\chapter{Breast Lump}

\section*{Definition}
A new palpable mass, nodule, or area of thickening in the breast or chest tissue. Many lumps are benign, but a new or persistent lump needs clinical assessment because examination alone cannot reliably determine its cause.

\section*{When to See a Doctor}

\begin{itemize}
 \item Any new breast lump, particularly in women over 30; lump with skin dimpling, nipple retraction, bloody discharge, or axillary lymphadenopathy
\end{itemize}

\section*{How It Develops}
Breast lumps have different causes, including normal or hormone-related tissue changes, fluid-filled cysts, benign growths, infection, injury, and cancer. Hormonal changes can affect some benign lumps and cancers. Breast cancers are a diverse group of diseases involving acquired or inherited molecular changes; no single gene or receptor abnormality explains every cancer.

\section*{Causes}
\begin{itemize}
 \item Fibroadenoma (most common benign solid breast mass in young women)
 \item Simple breast cyst (fluid-filled sac, common in premenopausal women)
 \item Fibrocystic changes (nodularity, tenderness with menstrual cycle)
 \item Breast cancer (invasive ductal carcinoma most common)
 \item Intraductal papilloma (nipple discharge, often bloody)
 \item Mastitis or breast abscess (infection, often during lactation)
 \item Lipoma or hamartoma (benign fatty tumors)
 \item Fat necrosis (post-traumatic or post-surgical)
 \item Phyllodes tumor (rare, can be benign or malignant)
 \item Gynecomastia (male breast enlargement from hormonal imbalance)
\end{itemize}

\section*{Related Symptoms}
\begin{itemize}
 \item Breast pain
 \item Nipple discharge
 \item Axillary lump
 \item Skin changes
 \item Swelling
\end{itemize}

\section*{Medical Evaluation}
\begin{itemize}
 \item Your doctor will perform a clinical breast examination including inspection for skin changes (dimpling, peau d'orange) and palpation in both sitting and supine positions.
 \item Imaging depends on age, pregnancy or lactation, examination findings, and local guidance. Ultrasound is usually the initial test for people younger than 30; diagnostic mammography or tomosynthesis, often with ultrasound, is usually used from age 40, while either may be appropriate from age 30 to 39.
 \item A core-needle biopsy or, in selected situations, fine-needle aspiration may be recommended when imaging or examination is suspicious, or when the clinical and imaging findings do not match.
 \item A breast specialist is considered for suspicious findings, persistent or enlarging lumps, or discordance between examination and imaging; a clearly benign, concordant finding may instead be monitored.
\end{itemize}

\section*{Treatment Options}
\begin{itemize}
 \item Simple breast cysts are aspirated if symptomatic; fibroadenomas may be observed if biopsy-confirmed and stable or excised if enlarging or symptomatic.
 \item Breast cancer treatment involves multidisciplinary care including surgery (lumpectomy or mastectomy), radiation, chemotherapy, hormonal therapy, and targeted therapy depending on subtype and stage.
 \item Mastitis is managed according to its cause and severity. During lactation, supportive care and continued comfortable milk removal may help; antibiotics may be needed for bacterial infection, and a collection or lack of improvement needs reassessment.
 \item Fat necrosis is managed conservatively with observation as it typically resolves spontaneously over several weeks to months.
\end{itemize}

\section*{Self-Care and Home Management}
\begin{itemize}
 \item Become familiar with the usual appearance and feel of your breasts or chest and report a new or persistent change.
 \item Note when the lump appeared and whether it changes with the menstrual cycle, but do not repeatedly press or massage it.
 \item Do not attempt to puncture, aspirate, or drain a lump at home.
 \item Schedule an appointment promptly rather than waiting for a lump to resolve on its own.
\end{itemize}

\section*{References}
\begin{thebibliography}{99}
\bibitem{breast_lump:ref1} Nelson HD, Pappas M, et al. ``Screening for breast cancer: a systematic review.'' \textit{JAMA}. 2016;315(4):391-405.
\bibitem{breast_lump:ref2} Salzman B, Fleegle S, et al. ``Common breast problems.'' \textit{Am Fam Physician}. 2012;86(4):343-349.
\bibitem{breast_lump:ref3} American College of Radiology. ``ACR BI-RADS Atlas, 5th ed.'' ACR, 2018.
\bibitem{breast_lump:ref4} Harris JR, Lippman ME, et al. \textit{Diseases of the Breast}, 5th ed. Wolters Kluwer, 2014.
\bibitem{breast_lump:ref5} UpToDate. ``Clinical manifestations and evaluation of breast lumps in adults.'' Wolters Kluwer, 2023.
\end{thebibliography}
